\chapter{公孫丑章句下}
\kaishi*

%公孫丑章句下
\section{孟子論天時地利人和}

孟子曰：「天時不如地利，地利不如人和。」\jiazhu[1]{天時謂時日支干五行王相孤虛之屬也。地利，險阻城池之固也。人和，得民心之所和樂也。}

「三里之城，七里之郭，環而攻之而不勝。夫環而攻之，必有得天時者矣；然而不勝者，是天時不如地利也。」\jiazhu[1]{環城圍之，必有得天時之善處者，然而城有不下，是不如地利。}

「城非不高也，池非不深也，兵革非不堅利也，米粟非不多也；委而去之，是地利不如人和也。」\jiazhu[1]{有堅強如此而破之走者，不得民心，民不爲守。衞懿公之民曰：君其使鶴戰。若是之類也。}

「故曰：域民不以封疆之界，固國不以山谿之險，威天下不以兵革之利。」\jiazhu[1]{域民，居民也。不以封疆之界禁之使民懷德也。不依險阻之固，恃仁惠也；不馮兵革之威，仗道德也。}

「得道者多助，失道者寡助。寡助之至，親戚畔之；多助之至，天下順之。以天下之所順，攻親戚之所畔，故君子有不戰，戰必勝矣。」\jiazhu[1]{得道之君，何嚮不平？君子之道，貴不戰耳；如其當戰，戰則勝矣。}\par \jiazhu[1]{章指言：民和爲貴，貴於天地，故曰得乎丘民爲天子也。}

\section{孟子將朝王}

孟子將朝王，王使人來曰：「寡人如就見者也，有寒疾，不可以風。朝將視朝，不識可使寡人得見乎？」\jiazhu[1]{孟子雖仕於齊，處師賔之位，以道見敬，或稱以病，未嘗趨朝而拜也。王欲見之，先朝使人往謂孟子云：寡人如就見者，若言就孟子之館相見也。有惡寒之病，不可見風。儻可來朝，欲力疾臨視朝，因得見孟子也。不知可使寡人得相見否。}

對曰：「不幸而有疾，不能造朝。」\jiazhu[1]{孟子不恱王之欲使朝，故稱有疾。}

明日，出弔於東郭氏。公孫丑曰：「昔者辭以病，今日弔，或者不可乎？」\jiazhu[1]{東郭氏，齊大夫家也。昔者，昨日也。丑以爲不可。}

曰：「昔者疾，今日愈，如之何不弔？」\jiazhu[1]{孟子言我昨日病，今日愈，我何爲不可以弔。}

王使人問疾，醫來。\jiazhu[1]{王以孟子實病，遣人將醫來，且問疾也。}

孟仲子對曰：「昔者有王命，有采薪之憂，不能造朝。今病小愈，趨造於朝，我不識能至否乎？」\jiazhu[1]{孟仲子，孟子之從昆弟，學於孟子者也。權辭以對如此。憂，病也。曲禮云：有負薪之憂。}

使數人要於路，曰：「請必無歸，而造於朝。」\jiazhu[1]{仲子使數人要告孟子，君命宜敬，當必造朝也。}

不得已而之景丑氏宿焉。\jiazhu[1]{孟子迫於仲子之言，不得已而心不欲至朝，因之其所知齊大夫景丑之家而宿焉。具以語景子。}

景子曰：「內則父子，外則君臣，人之大倫也。父子主恩，君臣主敬。丑見王之敬子也，未見所以敬王也。」\jiazhu[1]{景丑責孟子不敬，何義也。}

曰：「惡！是何言也？齊人無以仁義與王言者，豈以仁義爲不美也？其心曰：『是何足與言仁義也』云爾，則不敬莫大乎是。」\jiazhu[1]{曰惡者，深嗟歎云景子之責我何言乎？今人言謂王無知，不足與言仁義云爾。絕語之辭也。人之不敬，無大於是者也。}

「我非堯舜之道不敢以陳於王前，故齊人莫如我敬王也。」\jiazhu[1]{孟子言我每見王，常陳堯舜之道以勸勉王，齊人豈有如我敬王者邪。}

景子曰：「否，非此之謂也。禮曰：『父召無諾；君命召，不俟駕。』固將朝也，聞王命而遂不果，宜與夫禮若不相似然。」\jiazhu[1]{景子曰：非謂不陳堯舜之道，謂爲臣固自當朝也。今有王命而不果行。果，能也。禮，父召無諾而不至也；君命召，輦車就牧，不坐待駕。而夫子若是，事宜與夫禮若不相似然乎？愚竊惑焉。}

曰：「豈謂是與？曾子曰：『晉楚之富，不可及也；彼以其富，我以吾仁；彼以其爵，我以吾義，吾何慊乎哉？』夫豈不義而曾子言之？是或一道也。」\jiazhu[1]{孟子荅景丑云：我豈謂是君臣召呼之閒乎？謂王不禮賢下士，故道曾子之言，自以不慊晉楚之君。慊，少也。曾子豈嘗言不義之事邪？是或者自得道之一義，欲以喻王，猶晉楚，我猶曾子，我臣輕於王乎。}

「天下有達尊三：爵一，齒一，德一。朝廷莫如爵，鄉黨莫如齒，輔世長民莫如德。惡得有其一以慢其二哉？」\jiazhu[1]{三者，天下之所通尊也。孟子謂賢者、長者有德有齒，人君無德但有爵耳，故云何得以一慢二乎。}

「故將大有爲之君，必有所不召之臣；欲有謀焉，則就之。其尊德樂道，不如是，不足與有爲也。」\jiazhu[1]{言古之大聖大賢，有所興爲之君，必就大賢臣而謀事，不敢召也。王者師臣，霸者友臣也。}

「故湯之於伊尹，學焉而後臣之，故不勞而王；桓公之於管仲，學焉而後臣之，故不勞而霸。」\jiazhu[1]{言師臣者王。桓公能師臣，而管仲不勉之於王，故孟子於上章陳其義，譏其功烈之卑也。}

「今天下地醜德齊，莫能相尚，無他，好臣其所敎，而不好臣其所受敎。」\jiazhu[1]{醜，類也。言今天下人君，土地相類，德敎齊等，不能相絕者，無他，但好臣其所敎敕役使之才可驕者耳，不能好臣大賢可從受敎者。}

「湯之於伊尹，桓公之於管仲，則不敢召。管仲且猶不可召，而況不爲管仲者乎？」\jiazhu[1]{孟子自謂不爲管仲，故非齊王之召己，己是以不往也。}\par \jiazhu[1]{章指言：人君以尊德樂義爲賢，君子以守道不回爲志。}

\section{陳臻問受金}

陳臻問曰：「前日於齊，王餽兼金一百而不受；於宋，餽七十鎰而受；於薛，餽五十鎰而受。前日之不受是，則今日之受非也；今日之受是，則前日之不受非也。夫子必居一於此矣。」\jiazhu[1]{陳臻，孟子弟子。兼金，好金也，其價兼倍於常者，故謂之兼金。一百，百鎰也。古者以一鎰爲一金，鎰，二十兩。}

孟子曰：「皆是也。當在宋也，予將有遠行，行者必以贐；辭曰：『餽贐。』予何爲不受？」\jiazhu[1]{贐，送行者贈賄之禮也，時人謂之贐。}

「當在薛也，予有戒心；辭曰：『聞戒，故爲兵餽之。』予何爲不受？」\jiazhu[1]{戒，有戒備不虞之心也。時有惡人欲害孟子，孟子戒備。薛君曰：聞有戒，此金可鬻以作兵備，故餽之，我何爲不受也。}

「若於齊，則未有處也。無處而餽之，是貨之也。焉有君子而可以貨取乎？」\jiazhu[1]{我在齊時無事，於義未有所處也。義無所處而餽之，是以貨財取我，欲使我懷惠也。安有君子而以貨財見取乎？}\par \jiazhu[1]{章指言：取與之道，必得其禮，於其可也，雖少不辭；義之無處，兼金不顧。}

\section{孟子之平陸}

孟子之平陸，謂其大夫曰：「子之持戟之士，一日而三失伍，則去之否乎？」\jiazhu[1]{平陸，齊下邑也。大夫，治邑大夫也。持戟，戰士也。一日三失其行伍，則去之否乎？去之，殺之也。戎昭果毅。}

曰：「不待三。」\jiazhu[1]{大夫曰：一失之則行罰，不及待三失伍也。}

「然則子之失伍也亦多矣。凶年饑歲，子之民，老羸轉於溝壑，壯者散而之四方者，幾千人矣。」\jiazhu[1]{轉，轉尸於溝壑也。此則子之失伍也。}

曰：「此非距心之所得爲也。」\jiazhu[1]{距心，大夫名。曰：此乃齊王之大政，不肯賑窮，非我所得專爲也。}

曰：「今有受人之牛羊而爲之牧之者，則必爲之求牧與芻矣。求牧與芻而不得，則反諸其人乎？抑亦立而視其死與？」\jiazhu[1]{牧，牧地。以此喻距心不得自專，何不致爲臣而去乎？何爲立視民之死也。}

曰：「此則距心之罪也。」\jiazhu[1]{距心自知以不去位爲罪也。}

他日，見於王曰：「王之爲都者，臣知五人焉。知其罪者，惟孔距心。」爲王誦之。王曰：「此則寡人之罪也。」\jiazhu[1]{孔，姓也。爲都，治都也。邑有先君之宗廟曰都。誦，言也。爲王言所與孔距心語者也。王知本之在己，故受其罪。}\par \jiazhu[1]{章指言：人臣以道事君，否則奉身以退。詩云：彼君子兮，不素餐兮。言不尸其禄也。}

\section{孟子謂蚳鼃}

孟子謂蚳鼃曰：「子之辭靈丘而請士師，似也，爲其可以言也。今既數月矣，未可以言與？」\jiazhu[1]{蚳鼃，齊大夫。靈丘，齊下邑。士師，治獄官也。周禮士師曰：以五戒先後刑罰，母使罪麗於民。孟子見蚳鼃辭外邑大夫，請爲士師，知其欲近王，似諫正刑罰之不中者。數月而不言，故曰未可以言與？以感責之也。}

蚳鼃諫於王而不用，致爲臣而去。\jiazhu[1]{三諫不用，致仕而去。}

齊人曰：「所以爲蚳鼃則善矣；所以自爲，則吾不知也。」\jiazhu[1]{齊人論者，譏孟子爲蚳鼃諫，使之諫而去則善矣，不知自諫又不去，故曰我不見其自爲謀者。}

公都子以告。\jiazhu[1]{公都子，孟子弟子。以齊人語告孟子也。}

曰：「吾聞之也：有官守者，不得其職則去；有言責者，不得其言則去。我無官守，我無言責也，則吾進退，豈不綽綽然有餘裕哉？」\jiazhu[1]{官守，居官守職者。言責，獻言之責，諫爭之官也。孟子言人云云，今我居師賔之位，進退自由，豈不綽綽然舒緩有餘裕乎？綽、裕皆寬也。}\par \jiazhu[1]{章指言：執職者劣，藉道者優。是以臧武仲雨行而不息，段干木偃寢而式閭。}

\section{孟子爲卿於齊}

孟子爲卿於齊，出弔於滕，王使蓋大夫王驩爲輔行。王驩朝暮見，反齊、滕之路，未嘗與之言行事也。\jiazhu[1]{孟子嘗爲齊卿，出弔滕君。蓋，齊下邑也。王以治蓋之大夫王驩爲輔行。輔，副使也。王驩，齊之諂人，有寵於王，後爲右師。孟子不恱其爲人，雖與同使而行，未嘗與之言行事，不願與之相比也。}

公孫丑曰：「齊卿之位，不爲小矣；齊、滕之路，不爲近矣。反之而未嘗與言行事，何也？」\jiazhu[1]{丑怪孟子不與驩議行事也。}

曰：「夫既或治之，予何言哉？」\jiazhu[1]{既，已也；或，有也。孟子曰：夫人既自謂有治行事，我將復何言哉？言其專知自善，不知諮於人也。}\par \jiazhu[1]{章指言：道不合者不相與言。王驩之操與孟子殊，君子處時，危行言遜，故不尤之，但不與言。至於公行之喪，以禮爲解也。}

\section{孟子自齊葬於魯}

孟子自齊葬於魯，反於齊，止於嬴。充虞請曰：「前日不知虞之不肖，使虞敦匠事。嚴，虞不敢請。今願竊有請也：木若以美然。」\jiazhu[1]{孟子仕於齊，喪母，歸葬於魯。嬴，齊南邑。充虞，孟子弟子。敦匠，厚作棺也。事嚴，喪事急。木若以泰美然也。}

曰：「古者棺椁無度，中古棺七寸，椁稱之。自天子達於庶人，非直爲觀美也，然後盡於人心。」\jiazhu[1]{孟子言：古者棺椁薄厚無尺寸之度。中古謂周公制禮以來，棺厚七寸，椁薄於棺，厚薄相稱相得也。從天子至於庶人，厚薄皆然，但重累之數、牆翣之飾有異。非直爲人觀視之美好也，厚者難腐朽，然後能盡於人心所不忍也。謂一世之厚，孝子更去辟世，是爲人盡心也。過是以往，變化自其理也。}

「不得，不可以爲恱；無財，不可以爲恱；得之爲有財，古之人皆用之，吾何爲獨不然？」\jiazhu[1]{恱者，孝子之欲厚送親，得之則恱也。王制所禁不得用之，不可以恱心也。無財以供，則度而用之。禮，喪事不外求，不可稱貸而爲恱也。禮得用之，財足備之，古人皆用之，我何爲獨不然？然，如是也。}

「且比化者無使土親膚，於人心獨無恔乎？」\jiazhu[1]{恔，快也。棺椁敦厚，比親體之變化，且無令土親肌膚，於人子之心獨不快然無所恨也。}

「吾聞之：君子不以天下儉其親。」\jiazhu[1]{我聞君子之道，不以天下人所得用之物儉約於其親，言事親竭其力者也。}\par \jiazhu[1]{章指言：孝必盡心，匪禮之踰。論語曰：生事之以禮，死葬之以禮，可謂孝矣。}

\section{沈同以其私問伐燕}

沈同以其私問曰：「燕可伐與？」孟子曰：「可。子噲不得與人燕，子之不得受燕於子噲。」\jiazhu[1]{沈同，齊大臣，自以其私情問，非王命也，故曰私。子噲，燕王也。子之，燕相也。孟子曰可者，以子噲不以天子之命而擅以國與子之，子之亦不受天子之命而私受國於子噲，故曰其罪可伐。}

「有仕於此，而子恱之，不告於王而私與之吾子之禄爵；夫士也，亦無王命而私受之於子，則可乎？何以異於是？」\jiazhu[1]{子謂沈同也。孟子設此以譬燕王之罪。}

齊人伐燕。\jiazhu[1]{沈同以孟子言可，因歸勸其王伐燕。}

或問曰：「勸齊伐燕，有諸？」\jiazhu[1]{有人問孟子勸齊王伐燕有之。}

曰：「未也。沈同問『燕可伐與』，吾應之曰『可』，彼然而伐之也。」\jiazhu[1]{孟子曰：我未勸王也。同問可伐乎？吾曰可。彼然而伐之。}

「彼如曰：『孰可以伐之？』則將應之曰：『爲天吏，則可以伐之。』」\jiazhu[1]{彼如將問我曰：誰可以伐之？我將曰：爲天吏則可以伐之。天吏，天所使，謂王者得天意者也。彼不復孰可，便自往伐之。}

「今有殺人者，或問之曰：『人可殺與？』則將應之曰：『可。』彼如曰：『孰可以殺之？』則將應之曰：『爲士師，則可以殺之。』今以燕伐燕，何爲勸之哉？」\jiazhu[1]{今有殺人者，問此人可殺否？將應之曰可。爲士官主獄則可以殺之矣。言燕雖有罪，猶當王者誅之耳。譬如殺人者雖當死，士師乃得殺之耳。今齊國之政猶燕政也，不能相踰，又非天吏也，我何爲當勸齊伐燕乎？}\par \jiazhu[1]{章指言：誅不義者必須聖賢，禮樂征伐自天子出，王道之正也。}

\section{燕人畔}

燕人畔。王曰：「吾甚慚於孟子。」\jiazhu[1]{燕人畔，不肯歸齊。齊王聞孟子與沈同言爲未勸王，今竟不能有燕，故慚之。}

陳賈曰：「王無患焉。王自以爲與周公孰仁且智？」\jiazhu[1]{陳賈，齊大夫也。問王曰：自視何如周公仁智乎？欲爲王解孟子意，故曰王無患焉。}

王曰：「惡！是何言也？」\jiazhu[1]{王歎曰：是何言？言周公何可及也。}

曰：「周公使管叔監殷，管叔以殷畔。知而使之，是不仁也；不知而使之，是不智也。仁智，周公未之盡也，而況於王乎？賈請見而解之。」\jiazhu[1]{賈欲以此說孟子也。}

見孟子，問曰：「周公何人也？」\jiazhu[1]{賈問之也。}

曰：「古聖人也。」\jiazhu[1]{孟子曰：周公，古之聖人也。}

曰：「使管叔監殷，管叔以殷畔也，有諸？」\jiazhu[1]{賈問有之否乎。}

曰：「然。」\jiazhu[1]{孟子曰：如是也。}

曰：「周公知其將畔而使之與？」\jiazhu[1]{賈問之也。}

曰：「不知也。」\jiazhu[1]{孟子曰：周公不知其將畔。}

「然則聖人且有過與？」\jiazhu[1]{過，謬也。賈曰：聖人且猶有謬誤。}

曰：「周公，弟也；管叔，兄也。周公之過，不亦宜乎？」\jiazhu[1]{孟子以爲周公雖知管叔不賢，亦不必知其將畔。周公惟管叔弟也，故愛之；管叔念周公兄也，故望之。親親之恩也。周公於此過謬，不亦宜乎？}

「且古之君子，過則改之；今之君子，過則順之。古之君子，其過也，如日月之食，民皆見之；及其更也，民皆仰之。今之君子，豈徒順之，又從爲之辭。」\jiazhu[1]{古之所謂君子，眞聖人賢人君子也。周公雖有此過，乃誅三監，作大誥，明敕庶國，是周公改之也。今之所謂君子，非眞君子也，順過飾非，就爲之辭。孟子言此以譏賈不能匡君，而欲以辭解之。}\par \jiazhu[1]{章指言：聖人親親，不文其過；小人順非，以諂其上也。}

\section{孟子致爲臣而歸}

孟子致爲臣而歸。\jiazhu[1]{辭齊卿而歸其室也。}

王就見孟子，曰：「前日願見而不可得，得侍同朝，甚喜。今又棄寡人而歸，不識可以繼此而得見乎？」\jiazhu[1]{謂未來仕齊也，遙聞孟子之賢而不能得見之。來就爲卿，君臣同朝，得相見，故喜也。今致爲臣，棄寡人而歸也。不知可以續今日之後，還使寡人得相見否。}

對曰：「不敢請耳，固所願也。」\jiazhu[1]{孟子對王言：不敢自請耳，固心之所願也。孟子意欲使王繼今當自來謀也。}

他日，王謂時子曰：「我欲中國而授孟子室，養弟子以萬鍾，使諸大夫國人皆有所矜式。子盍爲我言之？」\jiazhu[1]{時子，齊臣也。王欲於國中央爲孟子築室，使養敎一國君臣之子弟，與之萬鍾之禄。中國者，使學者遠近鈞也。矜，敬也；式，法也。欲使諸大夫國人皆敬法其道。盍，何不也。謂時子何不爲我言之於孟子，知肯就之否。}

時子因陳子而以告孟子，陳子以時子之言告孟子。\jiazhu[1]{陳子，孟子弟子陳臻。}

孟子曰：「然，夫時子惡知其不可也？如使予欲富，辭十萬而受萬，是爲欲富乎？」\jiazhu[1]{孟子曰：如是夫時子安能知其不可乎？時子以我爲欲富，故以禄誘我。我往者享十萬鍾之禄，以大道不行故去耳。今更當受萬鍾，是爲欲富乎？距時子之言也。}

「季孫曰：『異哉子叔疑！』」\jiazhu[1]{二子，孟子弟子也。季孫知孟子意不欲，而心欲使孟子就之，故曰異哉弟子之所聞也。子叔心疑，亦以爲可就也。}

「使已爲政，不用，則亦已矣，又使其子弟爲卿。人亦孰不欲富貴？而獨於富貴之中有私龍斷焉。」\jiazhu[1]{孟子解二子之異意疑心，曰：齊王使我爲政，不用，則亦自止矣。今又欲以其子弟故，使我爲卿而與我萬鍾之禄。人亦誰不欲富貴乎？是猶獨於富貴之中有此私登龍斷之類也。我則恥之。}

「古之爲市也，以其所有易其所無者，有司者治之耳。有賤丈夫焉，必求龍斷而登之，以左右望而罔市利。人皆以爲賤，故從而征之。征商自此賤丈夫始矣。」\jiazhu[1]{古者市置有司，但治其爭訟，不征稅也。賤丈夫，貪人可賤者也。入市則求龍斷而登之。龍斷，謂堁斷而高者也。左右占望，見市中有利，罔羅而取之。人皆賤其貪也，故就征取其利。後世緣此遂征商人。孟子言我苟貪萬鍾，不恥屈道，亦與此賤丈夫何異也？古者謂周公以前，周禮有關市之賦也。}\par \jiazhu[1]{章指言：君子正身行道，道之不行，命也。不爲利回，創業可繼。是以君子以龍斷之人爲惡戒也。}

\section{孟子去齊宿於晝}

孟子去齊，宿於晝。有欲爲王留行者，\jiazhu[1]{晝，齊西南近邑也。孟子去齊欲歸鄒，至晝而宿也。齊人之知孟子者追送見之，欲爲王留孟子之行。}

坐而言。不應，隱几而臥。\jiazhu[1]{客危坐而言留孟子之言也。孟子不應荅，因隱倚其几而臥也。}

客不恱曰：「弟子齊宿而後敢言，夫子臥而不聽，請勿復敢見矣。」\jiazhu[1]{齊，敬；宿，素也。弟子素持敬心來言，夫子慢我，不受我言。言而遂起，退欲去，請絕也。}

曰：「坐！我明語子。昔者魯繆公無人乎子思之側，則不能安子思；泄柳、申詳無人乎繆公之側，則不能安其身。」\jiazhu[1]{往者魯繆公尊禮子思，子思以道不行則欲去。繆公常使賢人往留之，說以方且聽子爲政，然則子思復留。泄柳、申詳亦賢者也，繆公尊之不如子思。二子常有賢者在繆公之側，勸以復之，其身乃安也。}

「子爲長者慮，而不及子思。子絕長者乎？長者絕子乎？」\jiazhu[1]{長者，老者也。孟子年老，故自稱長者。言子爲我慮，不如子思時賢人也。不勸王使我得行道，而伹勸我留，留者何爲哉？此爲子絕我乎？又我絕子乎？何爲而愠恨也。}\par \jiazhu[1]{章指言：惟賢能安賢，智能知微，以愚喻智，道之所以乖也。}

\section{孟子去齊尹士語人}

孟子去齊，尹士語人曰：「不識王之不可以爲湯武，則是不明也；識其不可，然且至，則是干澤也。千里而見王，不遇故去，三宿而後出晝，是何濡滯也？士則玆不恱。」\jiazhu[1]{尹士，齊人也。干，求也；澤，禄也。尹士與論者言之，云孟子不知，則爲求禄。濡滯，猶稽也。既去近，留於晝三日，怪其猶久，故云士於此事不恱也。}

高子以告。\jiazhu[1]{高子亦齊人，孟子弟子。以尹士之言告孟子也。}

曰：「夫尹士惡知予哉？千里而見王，是予所欲也；不遇故去，豈予所欲哉？予不得已也。」\jiazhu[1]{孟子曰：夫尹士安能知我哉？我不得已而去耳，何汲汲而驅馳乎？}

「予三宿而出晝，於予心猶以爲速，王庶幾改之！王如改諸，則必反予。」\jiazhu[1]{我自謂行速疾矣，兾王庶幾能反覆招還我矣。}

「夫出晝，而王不予追也，予然後浩然有歸志。」\jiazhu[1]{浩然，心浩浩有遠志。}

「予雖然，豈舍王哉？王由足用爲善。王如用予，則豈徒齊民安，天下之民舉安。王庶幾改之，予日望之！」\jiazhu[1]{孟子以齊大國，知其可以行善政，故戀戀望王之改而反之，是以安行也。豈徒齊民安，言君子達則兼善天下也。}

「予豈若是小丈夫然哉？諫於其君而不受，則怒，悻悻然見於其面，去則窮日之力而後宿哉？」\jiazhu[1]{我豈若狷急小丈夫，恚怒其君而去，極日力而宿，懼其不遠者哉？論曰：悻悻然小人哉。言己志大在於濟一世之民，不爲小節也。}

尹士聞之，曰：「士誠小人也。」\jiazhu[1]{尹士聞義則服。}\par \jiazhu[1]{章指言：大德洋洋，介士察察。賢者志其大者，不賢者志其小者，此之謂也。}

\section{孟子去齊充虞路問}

孟子去齊，充虞路問曰：「夫子若有不豫色然。前日虞聞諸夫子曰：『君子不怨天，不尤人。』」\jiazhu[1]{路，道也。於路中問也。充虞謂孟子去齊有恨心，顔色不恱也。}

曰：「彼一時，此一時也。五百年必有王者興，其間必有名世者。由周而來，七百有餘歲矣。以其數，則過矣；以其時考之，則可矣。」\jiazhu[1]{彼前聖賢之出是有時也，今此時亦是其一時也。五百年有王者興，有興王道者也。名世，次聖之才，物來能名正一世者，生於聖人之閒也。七百有餘歲，謂周家王迹始興，大王、文王以來。考驗其時，則可有也。}

「夫天未欲平治天下也；如欲平治天下，當今之世，舍我其誰也？吾何爲不豫哉？」\jiazhu[1]{孟子自謂能當名世之士，時又值之而不得施，此乃天自未欲平治天下耳，非我之愆。我固不怨天，何爲不恱豫乎？}\par \jiazhu[1]{章指言：聖賢興作，與天消息。天非人不因，人非天不成。是故知命者不憂不懼也。}

\section{孟子去齊居休}

孟子去齊，居休。公孫丑問曰：「仕而不受禄，古之道乎？」\jiazhu[1]{休，地名。丑問古人之道仕不受禄邪？怪孟子於齊不受禄也。}

曰：「非也。於崇，吾得見王，退而有去志，不欲變，故不受也。」\jiazhu[1]{崇，齊地。孟子言不受禄，非古之道也。於崇吾始得見齊王，知其不能納善，退出志欲去矣。不欲即去，若爲變詭見非泰甚，故旦宿留心欲去，故不復受禄。}

「繼而有師命，不可以請。久於齊，非我志也。」\jiazhu[1]{言我本志欲速去，繼見之後有師旅之命，不得請去，故使我久而不受禄耳。久非我本志也。}\par \jiazhu[1]{章指言：禄以食功，志以率事。無其事而食其禄，君子不由也。}
